En conclusión, las pruebas nos demostraron que la notación teórica Big-O no cuenta toda la historia a la hora de predecir el rendimiento real. Factores físicos del hardware,
como aprovechar bien la memoria caché o el alto costo de estar pidiendo memoria dinámica, cambian drásticamente los tiempos de ejecución.
Esto lo vimos clarísimo cuando los algoritmos \textit{in-place} como Quick Sort y std::sort estuvieron por encima de Merge Sort, o cuando Naive
le ganó con creces a Strassen simplemente porque este último satura al procesador creando submatrices en cada paso recursivo.
En resumen, para desarrollar un buen software no basta con que la matemática del algoritmo sea eficiente, tenemos que programarlo pensando en que la máquina maneja sus recursos por dentro.